% Response to Reviewer 1
% Same text as rebuttal/reviewer_1.md (1479 characters as posted).
\section*{Response to Reviewer 1}

Thank you for your careful review and for appreciating our content-based design.

\textbf{1. Method:} No similarity threshold is used; $\tau$ on p.~3 is a leftover from an earlier version, sorry for the confusion. Algorithm 1's reuse fraction $\gamma$ instead fixes how many tokens are skipped within budget $k$. The memory holds the tokens the host recomputes in the first $W=4$ frames, merged by $T$ bipartite passes with update $d \leftarrow (d + \sum a)/(1 + n)$, where $a$ are the $n$ tokens matched to $d$. It is rebuilt every $N$ frames (not used on our $\sim$10 s clips). Later frames skip the $\gamma k$ proposed tokens most similar to the memory. We use $\gamma = 0.25$, $T=6$ for detection (0.95, 2 for K400), chosen on a 10\% subset (p.~4). The revision will add this.

\textbf{2. Numbers:} Thank you; we will make these numbers consistent throughout.

\textbf{3. Practicality:} For latency, PSM is 14\% faster than dense (60.0 vs.\ 70.1 ms) and close to the host. For memory, it adds $\sim$5 MB, independent of video length; the other 2,324 MB is the host's (p.~5). For accuracy, it costs 1.9 mAP vs.\ the host at 672 (3.9 at 1024) for Table 1's lowest compute (60.7 $\rightarrow$ 46.6 GFLOPs/frame). On the trade-off, we agree a budget-matched host recovers $\sim$1.5 mAP (p.~5), and we will narrow our claims. PSM is a drop-in, training-free layer ($\gamma = 0$ is the host) for recurring content, e.g.\ stationary cameras (pp.~4--5).

\textbf{4. Presentation:} Thank you; we will fix the Fig.~1 caption, if possible add the granularity table, and revise the headings.

Thank you again!
